% ============================================================
%  Nature Communications manuscript
%  "Waveform morphology of alpha oscillations predicts
%   network connectivity in human stereoelectroencephalography"
%
%  Compile with: pdflatex main.tex  (twice for cross-references)
%  Requires: booktabs, graphicx, amsmath, amssymb, natbib,
%            geometry, hyperref, multirow, xcolor, caption
% ============================================================

\documentclass[11pt, a4paper]{article}

% ── packages ─────────────────────────────────────────────────────────────────
\usepackage[margin=2.5cm]{geometry}
\usepackage{amsmath, amssymb}
\usepackage{graphicx}
\usepackage{booktabs}
\usepackage{multirow}
\usepackage{natbib}
\usepackage[colorlinks=true, linkcolor=blue, citecolor=blue, urlcolor=blue]{hyperref}
\usepackage{caption}
\usepackage{subcaption}
\usepackage{xcolor}
\usepackage{setspace}
\usepackage{lineno}
\usepackage{microtype}

\onehalfspacing
\linenumbers

% ── custom commands ───────────────────────────────────────────────────────────
\newcommand{\rho}{\ensuremath{\varrho}}  % use varrho for Spearman rho
\renewcommand{\rho}{\ensuremath{\rho}}
\newcommand{\imf}{\textsc{imf}}
\newcommand{\seeg}{\textsc{seeg}}
\newcommand{\car}{\textsc{car}}
\newcommand{\emd}{\textsc{emd}}
\newcommand{\imcoh}{$|$ImCoh$|$}
\newcommand{\pdc}{\textsc{pdc}}
\newcommand{\tra}{\textsc{tra}}
\newcommand{\var}{\textsc{var}}
\newcommand{\fdr}{\textsc{fdr}}
\newcommand{\lme}{\textsc{lme}}
\newcommand{\asc}{asc2desc}
\newcommand{\ptr}{peak2trough}

% ─────────────────────────────────────────────────────────────────────────────
\begin{document}

% ── title block ──────────────────────────────────────────────────────────────
\begin{titlepage}
\vspace*{2cm}
{\LARGE\bfseries
Waveform morphology of alpha oscillations predicts network
connectivity in human stereoelectroencephalography
\par}

\vspace{1.5cm}
{\large
Author~1$^{1}$,
Author~2$^{2}$,
Author~3$^{1,3}$
\par}

\vspace{0.5cm}
{\normalsize
$^{1}$Department of Neuroscience, [Institution], [City], [Country]\\
$^{2}$Department of Biomedical Engineering, [Institution], [City], [Country]\\
$^{3}$[Additional affiliation]
\par}

\vspace{1.5cm}
{\small\noindent
\textbf{Correspondence:} [Corresponding author name] (\texttt{email@institution.edu})
\par}

\vspace{1cm}
{\small\noindent
\textbf{Keywords:}
stereoelectroencephalography; alpha oscillations; waveform morphology;
empirical mode decomposition; brain connectivity; imaginary coherence;
partial directed coherence; electrical brain stimulation
\par}
\end{titlepage}

% ─────────────────────────────────────────────────────────────────────────────
\newpage
\section*{Abstract}

Neural oscillations are ubiquitous features of brain activity, yet the
canonical analysis treats them as sinusoidal and characterises them only by
power and frequency. Accumulating evidence demonstrates that the \emph{shape}
of oscillatory waveforms carries functionally relevant information beyond these
spectral properties. Here we tested whether waveform morphology of
intracranially recorded alpha oscillations predicts inter-electrode network
connectivity after electrical brain stimulation. We applied empirical mode
decomposition (EMD) to stereoelectroencephalography (\seeg) recordings from
eight participants, isolating the dominant alpha component (intrinsic mode
function~4; mean instantaneous frequency 9.9~$\pm$~0.8~Hz). Three cycle-level
shape descriptors were extracted: instantaneous frequency (IF), the
ascent-to-descent ratio (\asc), and the peak-to-trough ratio (\ptr). We
computed three complementary pairwise connectivity measures---imaginary coherence
(ImCoh), generalised partial directed coherence (gPDC), and temporal
irreversibility (TRA)---and aggregated them to the electrode level. Electrical
stimulation selectively prolonged the ascending phase of the alpha waveform at
stimulated electrodes (mean $\Delta$\asc~$= 0.016 \pm 0.020$;
Mann--Whitney $p = 5.1 \times 10^{-8}$), with no change at non-stimulated sites.
Pooled Spearman correlations revealed robust relationships between waveform
features and connectivity across more than 1,000 electrode observations. Linear
mixed-effects modelling confirmed that the ascent/descent ratio independently
predicted all three connectivity measures after controlling for brain region,
stimulation frequency, and channel status. These findings establish oscillatory
waveform shape as a novel and independent predictor of large-scale neural
network connectivity.

% ─────────────────────────────────────────────────────────────────────────────
\newpage
\section{Introduction}

Neural oscillations are among the most prominent and reproducible features of
electrophysiological brain activity, observed across species, recording
modalities, and cognitive states~\citep{buzsaki2004neuronal, jensen2007human,
fries2015rhythms}. The alpha band (8--12~Hz) in particular has attracted
sustained attention as a marker of cortical excitability, attentional
suppression, and long-range synchronisation~\citep{klimesch1999eeg,
jensen2002oscillations, palva2007new}. The prevailing analytical approach treats
alpha oscillations as quasi-sinusoidal, characterising them primarily by their
spectral power and peak frequency. This approach, however, discards substantial
structure embedded in the \emph{shape} of individual oscillatory cycles---a
structure increasingly recognised as physiologically and computationally
meaningful~\citep{cole2017brain, donoghue2022parameterizing}.

Any nonsinusoidal waveform feature in an oscillation produces within-cycle
modulations of the instantaneous frequency: the oscillation speeds up and
slows down as the waveform departs from a perfect sinusoid. This fundamental
property, formalised by Quinn et al.~\citeyearpar{quinn2021within}, makes the
instantaneous frequency profile an ideal readout of cycle shape---a cycle with
a slow ascending edge has low instantaneous frequency during the ascending phase;
a cycle with a wide, blunt peak has low instantaneous frequency near the peak.
By combining masked empirical mode decomposition (\emd) with phase-alignment---a
process that reregisters each cycle's instantaneous frequency profile onto a
uniform phase grid---Quinn et al.~\citeyearpar{quinn2021within} introduced an
analytical framework that isolates the dominant oscillatory component and
quantifies its waveform shape at the level of individual cycles, without
assumptions about stationarity or sinusoidality~\citep{huang1998empirical}.
Applied to mouse hippocampal theta oscillations, this approach revealed that
cycle shape varies substantially on a cycle-by-cycle basis and that specific
waveform motifs---indexed by the ascending-to-descending duration ratio and the
peak-to-trough duration ratio (``control point'' features)---relate to
behavioural and network variables. Whether similar waveform shape dynamics carry
information about network connectivity in the human brain remains an open
question.

This is an important gap to address for two reasons. First, connectivity measures
such as imaginary coherence (ImCoh)~\citep{nolte2004identifying} and generalised
partial directed coherence (gPDC)~\citep{baccala2001partial, baccala2007generalized}
index distinct aspects of network organisation---synchrony robustness to volume
conduction and directional Granger-causal influence, respectively---and if
waveform shape covaries with these measures, it implies that oscillatory
morphology reflects the network context of individual recording sites, not merely
local membrane dynamics. This would open a new route to characterising
functional connectivity from single-electrode recordings, without requiring
pairwise computations. Second, electrical brain stimulation systematically
perturbs neural circuit dynamics and may alter the very waveform properties
indexed by the Quinn framework~\citep{bestmann2015understanding,
voytek2014oscillatory, schaworonkow2019waveform}; establishing whether and how
stimulation-induced waveform changes track connectivity reorganisation would
provide a mechanistically grounded biomarker for therapeutic neuromodulation.

Here, we tested whether the shape of individual alpha oscillatory cycles
predicts network connectivity in intracranial stereoelectroencephalography
(\seeg) recordings from eight participants who received electrical brain
stimulation at 1~Hz and 50~Hz. We applied masked-\emd to decompose the
broadband \seeg signal at every electrode, isolated the alpha intrinsic mode
function, and derived three cycle-level waveform features---mean instantaneous
frequency, the ascending-to-descending phase duration ratio (\asc), and the
peak-to-trough duration ratio (\ptr)---following the framework of Quinn et
al.~\citeyearpar{quinn2021within}. We computed three complementary
electrode-level connectivity measures (ImCoh, gPDC asymmetry, and temporal
irreversibility; TRA) and tested whether waveform shape predicted connectivity
both across all electrodes and after controlling for anatomical and experimental
confounds in a linear mixed-effects model. We show that electrical stimulation
selectively prolongs the ascending phase of the alpha waveform at stimulated
sites, and that this \asc feature independently predicts all three connectivity
measures across brain regions and stimulation protocols---establishing
oscillatory waveform morphology as a systematic predictor of large-scale
network connectivity in the human brain.

% ─────────────────────────────────────────────────────────────────────────────
\section{Results}

\subsection{EMD isolates a stable alpha component across all electrodes}

We applied masked empirical mode decomposition (masked-\emd;
\citealt{quinn2021within}) to the post-stimulation epoch of each electrode
independently, decomposing the broadband \seeg signal into six intrinsic mode
functions (\imf{}s). Across all 1,009 electrodes and eight participants,
\imf{}~4 (zero-indexed index~3) consistently captured the dominant oscillation
in the alpha frequency range (mean instantaneous frequency:
$9.9 \pm 0.8$~Hz; range 6.4--11.9~Hz; Table~\ref{tab:demographics}), in
agreement with prior reports using the same decomposition
framework~\citep{quinn2021within}. The decomposition procedure, the
extraction of three cycle-level waveform features (instantaneous frequency,
ascent/descent ratio, and peak/trough ratio), and the resulting population
distributions are illustrated in Figure~\ref{fig:fig1}.

Three features were derived from each accepted alpha cycle
(see Methods): the mean instantaneous frequency of IMF~4 within the cycle
($\overline{\text{IF}}$; mean across electrodes $9.88 \pm 0.76$~Hz); the
ascent/descent ratio (\asc; mean $0.506 \pm 0.006$), quantifying the fraction
of the cycle spent in the ascending phase; and the peak/trough ratio (\ptr;
mean $0.490 \pm 0.003$), indexing the position of the descending zero-crossing
within the cycle. Electrode-level values were obtained by averaging each metric
across all accepted cycles within the epoch. The resulting distributions were
unimodal and well-separated from chance, confirming that \emd successfully
isolated individual alpha oscillations rather than broadband noise.

\subsection{Electrical stimulation selectively alters ascending-phase duration
at stimulated electrodes}

We examined whether electrical stimulation changed the morphology of the alpha
waveform by comparing pre- and post-stimulation \asc and \ptr values at
stimulated ($n = 42$) and non-stimulated ($n = 1{,}248$) electrodes
(Figure~\ref{fig:fig2}). The grand-mean normalised waveform at stimulated
electrodes shifted markedly between epochs: the post-stimulation waveform
exhibited a prolonged ascending phase relative to baseline, whereas the overall
waveform envelope and amplitude were preserved. In contrast, non-stimulated
electrodes showed negligible pre-to-post differences, with the mean waveforms
from both epochs overlapping closely across the entire cycle.

Quantitatively, stimulated electrodes showed a significantly larger
pre-to-post increase in \asc than non-stimulated electrodes
($\Delta$\asc: stimulated $0.016 \pm 0.020$ vs non-stimulated
$-0.001 \pm 0.011$; Mann--Whitney $U$, $p = 5.1 \times 10^{-8}$), whereas
$\Delta$\ptr did not differ between groups ($p = 0.17$; Figure~\ref{fig:fig2},
panels C--D). The specificity of the effect to \asc, and to stimulated sites,
indicates that the stimulation selectively extended the ascending phase of the
alpha cycle without altering peak sharpness, consistent with a selective
modulation of the excitatory drive underlying the ascending phase of the
oscillation.

\subsection{Waveform shape predicts network connectivity across electrodes}

We next tested whether the three waveform features predicted the three
electrode-level connectivity measures (ImCoh, PDC~asymmetry, and TRA) across
the full post-stimulation dataset ($n = 1{,}009$ electrodes; Figure~\ref{fig:fig3}).
Pooled Spearman correlations across all electrodes revealed several significant
associations. Instantaneous frequency was the strongest raw predictor of
directed and irreversible connectivity: higher IF was associated with lower PDC
asymmetry ($\rho = -0.158$, $p < 0.001$) and lower TRA ($\rho = -0.207$,
$p < 0.001$), consistent with a pattern in which higher-frequency alpha
oscillations appear at electrodes that are net receivers (rather than net
drivers) in the network, and at sites with lower temporal irreversibility.
ImCoh showed a positive association with IF ($\rho = 0.075$, $p = 0.018$).
The ascent/descent ratio predicted PDC asymmetry significantly
($\rho = -0.086$, $p = 0.007$), and peak/trough ratio showed a significant
association with ImCoh ($\rho = -0.077$, $p = 0.014$).

We additionally tested the within-subject consistency of these relationships by
computing Spearman $\rho$ separately within each of the eight participants and
applying the Wilcoxon signed-rank test across the resulting distribution of
$\rho$ values. PDC asymmetry~$\sim$~\asc reached significance under this
stricter, subject-wise test (median $\rho = -0.085$; Wilcoxon $p = 0.016$),
while TRA~$\sim$~IF showed a trend (median $\rho = -0.109$; Wilcoxon
$p = 0.078$). Consistent directions of effect were observed across all eight
participants for the strongest associations, indicating that the pooled
relationships were not driven by individual outlier subjects.

\subsection{Confound-controlled LME models confirm independent waveform
shape effects}

The pooled Spearman correlations shown above do not control for potential
confounds including anatomical brain region, stimulation frequency protocol,
and stimulated-channel status---all of which could independently influence
both waveform shape and connectivity. To address this, we fitted separate
linear mixed-effects (\lme) models for each connectivity outcome, regressing
the z-scored connectivity measure onto all three z-scored waveform features
simultaneously, with fixed-effect covariates for brain region, stimulation
frequency, and channel status, and a random intercept per participant
(Figure~\ref{fig:fig4}; Table~\ref{tab:lme}).

The \lme analysis revealed that the ascent/descent ratio (\asc) was a
significant, independent predictor of all three connectivity measures after
controlling for confounds (Figure~\ref{fig:fig5}). Higher \asc predicted
greater ImCoh ($\hat{\beta} = 0.129$, 95\%~CI [0.073, 0.185],
$q < 0.001$), lower PDC~asymmetry ($\hat{\beta} = -0.097$, 95\%~CI
[$-$0.158, $-$0.035], $q = 0.004$), and greater TRA
($\hat{\beta} = 0.094$, 95\%~CI [0.039, 0.149], $q = 0.003$).
Instantaneous frequency independently predicted ImCoh ($\hat{\beta} = 0.117$,
95\%~CI [0.046, 0.187], $q = 0.003$), while peak/trough ratio was the
strongest predictor of PDC asymmetry ($\hat{\beta} = 0.168$,
95\%~CI [0.113, 0.224], $q < 0.001$). Together, five of the nine
waveform-feature--connectivity combinations survived \fdr-correction
(Benjamini--Hochberg; $q < 0.05$), with \asc being the only predictor
significant across all three connectivity outcomes (Table~\ref{tab:lme}).

The fact that the \asc effect on ImCoh, TRA, and PDC emerged robustly
after controlling for brain region argues against the possibility that the
observed associations simply reflect regionally specific patterns of oscillatory
morphology and connectivity. Furthermore, the inclusion of stimulated-channel
status as a covariate demonstrates that the waveform--connectivity relationship
is not merely a consequence of stimulation altering both variables
simultaneously: the effect persists even within the non-stimulated electrode
population, where waveform shape varies naturally in the absence of direct
perturbation.

\subsection{Waveform--connectivity associations are consistent across
stimulation protocols and brain regions}

To assess the robustness of the main findings across experimental and
anatomical subgroups, we examined IF~vs~connectivity separately for stimulated
and non-stimulated electrodes and for 1~Hz and 50~Hz stimulation protocols
(Figure~\ref{fig:fig6}). Non-stimulated electrodes ($n = 981$) showed
consistent and statistically robust associations between IF and all three
connectivity measures (ImCoh: $\rho = 0.074$, $p = 0.020$;
PDC: $\rho = -0.194$, $p < 0.001$; TRA: $\rho = -0.171$, $p < 0.001$).
The equivalent correlations in stimulated electrodes ($n = 28$) were not
significant, likely reflecting the small number of stimulated sites.
Both the 1~Hz and 50~Hz protocols yielded significant IF~vs~PDC and
IF~vs~TRA associations, confirming that the relationships were not specific to
one stimulation paradigm.

Regionally, the strongest and most consistent associations were observed in
frontal ($n = 334$) and insula ($n = 110$) cortices (Figure~\ref{fig:fig7}).
Occipital electrodes showed a significant positive PDC association
(FDR-corrected $p < 0.05$), and insula showed significant TRA associations,
while hippocampal and parietal regions showed non-significant trends. The
regional heterogeneity suggests that the waveform--connectivity link is not a
uniform whole-brain phenomenon but may reflect the specific role of alpha
oscillations in different functional networks.

% ─────────────────────────────────────────────────────────────────────────────
\section{Discussion}

In this study we combined empirical mode decomposition of intracranial alpha
oscillations with three complementary measures of network connectivity to test
whether the shape of individual oscillatory waveforms predicts the functional
network context of recording electrodes. Our central finding is that the
ascent/descent ratio (\asc) of the alpha cycle---the fraction of the cycle
duration spent in the ascending phase---independently and consistently predicted
imaginary coherence, directional drive asymmetry, and temporal irreversibility
across more than 1,000 electrode observations spanning eight participants, six
brain regions, and two stimulation frequency protocols. This result is robust to
control for anatomical region, stimulation condition, and between-participant
variability, and holds both after direct electrical perturbation of neural
circuits and in the resting-state context of unperturbed electrodes.

\paragraph{Waveform shape as a connectivity biomarker.}
A key contribution of the present work is the demonstration that cycle-level
alpha waveform morphology---specifically the ascending-to-descending duration
ratio---tracks the network position of each recording electrode. Following
Quinn et al.~\citeyearpar{quinn2021within}, the \asc feature corresponds to
the ratio of time the oscillation spends rising versus falling, with values
above 0.5 indicating a slow, prolonged ascending phase. We show that
electrodes with larger \asc (slower ascending, faster descending) exhibited
higher imaginary coherence and greater temporal irreversibility, but lower PDC
outflow asymmetry. This pattern is interpretable within a framework in which
the ascending phase of the alpha cycle reflects the summation of convergent
excitatory synaptic input driving the membrane potential toward the oscillatory
peak~\citep{cole2017brain, schaworonkow2019waveform}. Recording sites that
receive strong afferent drive from distributed network nodes are expected to
show slower, broader ascending phases---consistent with higher \asc values---
while also appearing as net \emph{receivers} of directional drive (lower PDC
asymmetry). The significant negative \lme coefficient for \asc on PDC
asymmetry ($\hat{\beta} = -0.097$, $q = 0.004$) is thus mechanistically
coherent: a prolonged ascending phase at a given electrode is associated with
reduced net outgoing directional influence, positioning that site as a network
integrator rather than an initiator. In the instantaneous frequency domain
\citep{quinn2021within}, this corresponds to a period of low IF during the
ascending phase that is detectable even in single cycles, suggesting that
cycle-by-cycle variations in network drive are continuously reflected in the
waveform morphology of the recorded signal.

\paragraph{Instantaneous frequency and network drive.}
Instantaneous frequency emerged as a robust predictor of directed and
irreversible connectivity in the pooled and \lme analyses, particularly for
PDC asymmetry and TRA. Higher-frequency alpha oscillations were associated
with more strongly negative PDC~asymmetry (net receivers) and lower TRA,
suggesting that electrodes oscillating faster than average within the alpha
band are embedded in network positions characterised by strong incoming
drive and relatively time-symmetric dynamics. This is consistent with
entrainment models in which the oscillatory frequency of a downstream node
is partly determined by the driving frequency of its upstream inputs:
electrodes receiving strong afferent drive may oscillate at frequencies
slightly above the autonomous population rhythm, producing a
systematically higher instantaneous frequency in the recorded \imf
\citep{fries2015rhythms, jensen2007human}. The \lme peak/trough ratio
effect on PDC asymmetry ($\hat{\beta} = 0.168$, the largest standardised
coefficient in the analysis) further underscores that cycle asymmetry
across both the time and amplitude dimensions carries complementary
information about network architecture.

\paragraph{Specificity of stimulation effects.}
The pre- to post-stimulation comparison (Figure~\ref{fig:fig2}) revealed a
selective prolongation of the ascending phase at stimulated electrodes that
was absent at non-stimulated sites. This specificity is important for two
reasons. First, it demonstrates that the observed waveform changes are not
simply a consequence of time-on-task or state drift, but rather reflect a
genuine perturbation of local circuit dynamics by the stimulation current.
Second, the change is restricted to \asc and does not extend to \ptr,
suggesting that the stimulation specifically affected the excitatory input
dynamics underlying the ascending phase without broadly altering the
peak--trough shape of the waveform. Whether this morphological change
reflects a direct biophysical consequence of the stimulation current on
membrane conductances, or an indirect network-mediated effect, remains an
open question that could be addressed with computational modelling.

\paragraph{Limitations and future directions.}
Several limitations merit consideration. The sample comprised eight
participants undergoing intracranial monitoring for epilepsy surgery, and
although the use of common-average referencing, pre-processing noise
rejection, and subject-level random effects mitigates several confounds,
generalisation to healthy populations and other pathologies requires
further investigation. The electrode counts per participant ranged from 16
to 258, reflecting the clinical implantation coverage rather than a
systematic anatomical sampling plan, and the relatively sparse hippocampal
and parietal coverage limits inferences in those regions. The Wilcoxon
signed-rank test with $n = 8$ participants has limited statistical power,
and several effects that appear numerically consistent across subjects did
not survive the conservative $p < 0.05$ threshold at this level of
analysis; larger datasets are needed to establish the full scope of
within-subject waveform--connectivity associations.

Additionally, all connectivity measures were computed on the alpha \imf
signal rather than the broadband signal, which isolates the oscillatory
component of interest but means that broadband aperiodic activity
\citep{donoghue2022parameterizing} did not contribute to the connectivity
estimates. Whether similar waveform--connectivity relationships exist in
other frequency bands or in non-oscillatory signal components is an
important open question. Finally, the correlational design of the
connectivity analyses cannot establish causality: the association between
waveform shape and connectivity could reflect shared upstream drivers
rather than a mechanistic link. Future work using closed-loop stimulation
and perturbation paradigms could distinguish these accounts.

\paragraph{Conclusion.}
Our results establish that the morphology of intracranial alpha oscillations
carries systematic information about network connectivity that goes beyond
what is captured by oscillatory power and peak frequency alone. The
ascent/descent ratio emerges as a particularly robust and interpretable
predictor, linking the temporal fine structure of individual cycles to the
position of a recording site within the large-scale network. These findings
motivate the inclusion of waveform shape analysis in future studies of
connectivity, stimulation, and brain state, and suggest that oscillatory
morphology may serve as a biomarker of network organisation that is
accessible from single-electrode recordings.

% ─────────────────────────────────────────────────────────────────────────────
\section{Methods}

\subsection{Participants and recordings}

Stereo-EEG (\seeg) recordings were obtained from eight participants
(Subject01, Subject04, Subject07, HEJ\_Subject12--16) undergoing
pre-surgical intracranial monitoring for drug-resistant focal epilepsy.
Electrodes were implanted according to clinical requirements and their
anatomical positions were assigned to brain regions using the Automated
Anatomical Labelling (AAL) atlas~\citep{tzourio2002automated}. Each
participant received transcranial electrical stimulation at two
frequencies: 1~Hz and 50~Hz. Recordings from one session
(Subject04, 50~Hz, run~2) were excluded owing to data quality issues
identified during pre-processing. All procedures were approved by the
relevant institutional ethics committee, and written informed consent was
obtained from each participant.

\subsection{Signal referencing}

All analyses were performed using a \textbf{common-average reference (\car)}
scheme, in which each channel's signal was re-referenced by subtracting the
instantaneous mean across all simultaneously recorded channels. A minimum of
five channels was required for a valid \car computation. The \car was chosen
for three reasons. First, it attenuates signals that are common to all
electrodes---such as volume-conducted far-field activity and shared
amplifier noise---while preserving local oscillatory signals that differ
across channels~\citep{nunez2006electric}. Second, it is the standard
reference in \seeg analysis, where the wide spatial separation of contacts
and the restricted field spread of intracranial signals limit the risk of
reference contamination that affects scalp \textsc{eeg}. Third, the \car
preserves the sign and magnitude of the imaginary coherence (ImCoh), which
depends on the phase relationship between electrode pairs; referencing
schemes that introduce artificial correlations (such as linked mastoids)
would distort ImCoh estimates. With a minimum of five simultaneously
recorded channels, the \car is well-conditioned and introduces only a small
residual contribution from each channel into all others
($\leq 1/(N-1)$ per channel).

\subsection{Epoch definition}

Two analysis epochs were defined relative to the stimulation period for
each recording. The \textbf{pre-stimulation epoch} spanned $-30$~s to
$-2$~s before stimulation onset (28~s total; the 2~s buffer excluded
anticipatory artefacts). The \textbf{post-stimulation epoch} began 1~s
after stimulation offset (artefact buffer) and extended for 28~s:
31--59~s for the 1~Hz protocol (stimulation offset at 30~s) and
6--34~s for the 50~Hz protocol (stimulation offset at 5~s). Primary
results are reported for the post-stimulation epoch.

\subsection{Pre-processing and noise rejection}

Each epoch was demeaned and notch-filtered at 60, 120, and 180~Hz using
second-order IIR notch filters ($Q = 30$) to suppress mains interference.
Linear drift was removed when the slope-to-interquartile-range ratio
exceeded 0.5. Channels were excluded if they contained non-finite values,
exhibited peak-to-peak amplitude exceeding $50 \times$ the
interquartile range, or if their power in a 1~Hz band around the line
frequency exceeded a flat-spectrum threshold. These criteria were applied
independently to each epoch; a channel could be retained in one epoch and
excluded from another.

\subsection{Empirical mode decomposition and IMF selection}

\paragraph{Masked-SIFT decomposition.}
Each preprocessed epoch was decomposed into intrinsic mode functions
(\imf{}s) using the \textbf{masked} variant of the \textsc{sift} algorithm
(\texttt{emd.sift.mask\_sift}; \citealt{quinn2021within}), implemented in the
\texttt{emd} Python package. Masked \emd is preferred over standard \emd
because it explicitly separates components whose frequencies are close to
one another by interleaving the sifting iterations with sinusoidal masking
signals, preventing mode mixing. The masking signal at iteration $i$ is
defined as:
%
\begin{equation}
  m_i(t) = A_{\text{mask}} \cdot \sin(2\pi f_{\text{mask},i} \cdot t),
  \label{eq:mask}
\end{equation}
%
where $A_{\text{mask}}$ is set proportional to the signal amplitude
(mode \texttt{ratio\_sig}, ratio~1), and the masking frequencies
$f_{\text{mask},i} \in \{80, 40, 20, 10, 5, 2, 1\}$~Hz were chosen
to successively isolate components from high to low frequency. A
maximum of 6~\imf{}s were extracted per electrode epoch; sifting was
terminated by the Rilling stopping criterion (standard-deviation
threshold $\sigma = 0.05$), and envelope interpolation used
monotone Piecewise Cubic Hermite Interpolating Polynomial (PCHIP) splines.

\paragraph{IMF 4 as the alpha component.}
\imf{}~4 (zero-indexed) was selected as the alpha component for all
analyses. This follows the general convention of the masked-\emd framework
\citep{quinn2021within}, in which higher-order \imf{}s capture progressively
lower-frequency oscillatory components. Across all electrodes, \imf{}~4
captured a mean instantaneous frequency of $9.9 \pm 0.8$~Hz
(range 6.4--11.9~Hz), confirming its correspondence to the alpha band.

\subsection{Instantaneous phase, frequency, and amplitude}

The Hilbert transform was applied to \imf{}~4 to obtain the analytic
signal:
%
\begin{equation}
  z_4(t) = \text{IMF}_4(t) + i \cdot \mathcal{H}\{\text{IMF}_4(t)\},
\end{equation}
%
from which the \textbf{instantaneous phase} $\phi(t) = \angle z_4(t) \in [0, 2\pi)$,
\textbf{instantaneous amplitude} $A(t) = |z_4(t)|$, and
\textbf{instantaneous frequency}
%
\begin{equation}
  \text{IF}(t) = \frac{1}{2\pi} \frac{d\phi}{dt}
  \label{eq:if}
\end{equation}
%
were computed. The phase derivative was smoothed with a 7-sample (half-width
3~samples) sliding-window average (\texttt{smooth\_phase=3}) to suppress
numerical noise from discrete differentiation.
Cycle boundaries were identified as the ascending zero-crossings of $\phi$,
detected as sample indices where $\phi(t+1) - \phi(t) < -\pi$ (i.e.\
wrap-around from $\approx 2\pi$ back to $\approx 0$).

\subsection{Cycle-level waveform feature extraction}

\paragraph{Cycle quality control.}
Cycles were retained only if they satisfied all of the following criteria:
(i) completeness (\texttt{is\_good}~$= 1$, as assessed by
\texttt{emd.cycles.Cycles}); (ii) maximum instantaneous amplitude exceeding
the 10th percentile of the epoch's amplitude distribution; (iii) maximum IF
within the cycle $< 20$~Hz; (iv) IF range within the cycle $< 20$~Hz; and
(v) the electrode contributed at least five accepted cycles in the epoch.

\paragraph{Instantaneous frequency per cycle.}
The mean IF was computed by averaging Equation~\ref{eq:if} over all
time-points within the accepted cycle:
%
\begin{equation}
  \overline{\text{IF}}_k = \frac{1}{T_k} \sum_{t \in C_k} \text{IF}(t),
\end{equation}
%
where $T_k$ is the duration of cycle $k$ in samples and $C_k$ is the
set of sample indices belonging to that cycle.

\paragraph{Ascending-to-descending duration ratio (asc2desc).}
Following the control-point framework of Quinn et al.~\citeyearpar{quinn2021within},
we quantified waveform asymmetry using the \textit{ascending-to-descending
duration ratio}: the fraction of the cycle's total duration spent in the
ascending phase. Because each cycle is bounded by ascending zero-crossings
(phase wrapping from $\approx 2\pi$ back to $\approx 0$), the ascending
phase comprises two non-contiguous intervals---from cycle onset to the
peak, and from the trough to cycle end:
%
\begin{equation}
  \text{asc2desc}_k = \frac{t_{\text{peak},k} + (T_k - t_{\text{trough},k})}{T_k},
  \label{eq:asc2desc}
\end{equation}
%
where $t_{\text{peak},k}$ is the peak sample (detected by parabolic
interpolation via \texttt{emd.cycles.cf\_peak\_sample}) and
$t_{\text{trough},k}$ is the trough sample
(\texttt{emd.cycles.cf\_trough\_sample}), both measured from the start of
cycle $k$, and $T_k$ is cycle length in samples. A sinusoidal cycle has
\asc~$= 0.5$; values above 0.5 indicate a slow-ascending, fast-descending
waveform; values below 0.5 indicate the reverse. This asymmetry is
mechanistically informative: a longer ascending phase corresponds to a
sustained, slow rise in membrane potential driven by excitatory synaptic
currents, whereas a sharp descent reflects rapid inhibitory or
hyperpolarising dynamics~\citep{cole2017brain, schaworonkow2019waveform}.
In the instantaneous frequency domain, a slow ascending edge produces
\emph{low} IF during the ascending phase~\citep{quinn2021within}.

\paragraph{Peak-to-trough duration ratio (peak2trough).}
As a second control-point feature, we computed the position of the
descending zero-crossing as a fraction of total cycle length, indexing
the relative duration of the peak segment versus the trough segment:
%
\begin{equation}
  \text{peak2trough}_k = \frac{t_{\text{desc\_zero},k}}{T_k},
  \label{eq:peak2trough}
\end{equation}
%
where $t_{\text{desc\_zero},k}$ is the sample index of the descending
zero-crossing (\texttt{emd.cycles.cf\_descending\_zero\_sample}).
Values near 0.5 indicate a symmetric waveform with equal peak and trough
durations; values below 0.5 indicate a short peak (early descending
zero-crossing, narrow peak); values above 0.5 indicate a broad, prolonged
peak. In the instantaneous frequency representation, a wide peak corresponds
to low IF near the peak phase ($\approx \pi/2$) and a narrow peak to high IF
at the same phase~\citep{quinn2021within}. This feature is therefore
complementary to \asc, capturing peak--trough asymmetry rather than
ascending--descending asymmetry.

\paragraph{Electrode-level aggregation.}
Per-cycle values of $\overline{\text{IF}}_k$, \asc$_k$, and \ptr$_k$ were
averaged across all accepted cycles within the epoch to yield a single
scalar per electrode:
%
\begin{equation}
  X_e = \frac{1}{K_e} \sum_{k=1}^{K_e} X_k,
\end{equation}
%
where $K_e$ is the number of accepted cycles for electrode $e$. This
averaging scheme gives each accepted cycle equal weight within an electrode.

\subsection{Connectivity measures}

Three pairwise connectivity measures were computed for every eligible electrode
pair within each epoch, then aggregated to the electrode level.

\subsubsection{Imaginary Coherence (ImCoh)}

The imaginary part of complex coherence is insensitive to instantaneous
zero-phase synchrony caused by volume conduction~\citep{nolte2004identifying}.
For signals $x(t)$ and $y(t)$, the cross-spectral density was estimated
using the multitaper method via \texttt{mne-connectivity} on seven
consecutive non-overlapping 4-s segments spanning the full 28-s epoch.
For each segment $j$, the cross-spectral density at frequency $f$ is
$S_{xy}^{(j)}(f) = X^{(j)*}(f) Y^{(j)}(f)$, where $X^{(j)}$ and
$Y^{(j)}$ are the Fourier transforms. Coherence was defined as:
%
\begin{equation}
  C_{xy}(f) = \frac{\langle S_{xy}(f) \rangle}
                   {\sqrt{\langle S_{xx}(f) \rangle \langle S_{yy}(f) \rangle}},
\end{equation}
%
and imaginary coherence as $\text{ImCoh}(f) = \text{Im}[C_{xy}(f)]$.
The mean absolute imaginary coherence was computed by averaging $|\text{ImCoh}(f)|$
across the 8--11~Hz band. Electrode-level \imcoh{} was defined as the mean
$|\text{ImCoh}|$ across all pairs involving that electrode.

\subsubsection{Generalised Partial Directed Coherence (gPDC)}

Partial directed coherence~\citep{baccala2001partial} quantifies
directional Granger-causal influence in the frequency domain, derived from
the transfer function of a fitted multivariate autoregressive (VAR) model.
We used the \textit{generalised} formulation~\citep{baccala2007generalized},
which normalises by the noise covariance and is therefore scale-invariant.
For each electrode pair $(\text{A}, \text{B})$:

\begin{enumerate}
  \item The IMF~4 signals from both channels were downsampled to 128~Hz
        using polyphase resampling (\texttt{scipy.signal.resample\_poly}).
  \item A bivariate VAR model of order $p$ was fitted with order selection
        by the Akaike Information Criterion (AIC), constrained to
        $p \leq \min(30,\, N/20)$ where $N$ is the number of samples.
  \item VAR stability was verified by confirming that all eigenvalues of
        the companion matrix lay within the unit circle.
  \item Residual whiteness was assessed using the Ljung--Box test at lag
        $L = \min(20,\, N/5)$; pairs with $p_{\text{LB}} < 0.01$ were
        excluded as insufficiently white.
  \item For stable, white-residual VAR models, the gPDC spectrum was
        evaluated at 64 frequency points within 8--11~Hz. Let $\mathbf{A}(f)$
        denote the VAR transfer matrix in the frequency domain; the gPDC
        from channel~$j$ to channel~$i$ is:
        \begin{equation}
          |\bar{\pi}_{ij}(f)|^2 =
          \frac{|A_{ij}(f)|^2 / \sigma_{jj}}
               {\sum_k |A_{kj}(f)|^2 / \sigma_{jj}},
          \label{eq:gpdc}
        \end{equation}
        where $\sigma_{jj}$ is the $j$-th diagonal element of the VAR
        noise covariance matrix.
  \item Band-averaged gPDC values in 8--11~Hz were computed:
        $\overline{\text{gPDC}}_{\text{A}\to\text{B}}$ and
        $\overline{\text{gPDC}}_{\text{B}\to\text{A}}$.
\end{enumerate}

Electrode-level PDC~asymmetry was defined as:
%
\begin{equation}
  \text{PDC}_e = \bar{\text{gPDC}}_{\text{out},e} - \bar{\text{gPDC}}_{\text{in},e},
  \label{eq:pdc_asym}
\end{equation}
%
where ``out'' and ``in'' refer to averages over all valid pairs in which
electrode $e$ is the source or the target, respectively. Positive values
indicate net outgoing drive.

\subsubsection{Temporal Irreversibility (TRA)}

Temporal irreversibility quantifies the degree to which the cross-correlation
between the Hilbert envelopes of two signals is asymmetric under time reversal,
providing an unsigned measure of temporal non-stationarity in the coupling
\citep{zanin2012time}. For each pair:

\begin{enumerate}
  \item The Hilbert envelope of IMF~4 was computed for each channel and
        normalised by its standard deviation.
  \item The cross-correlation $[y \star x]_\ell$ was computed via FFT
        at lags $\ell = 1, \ldots, L$ where $L = \lfloor F_s \times 1 \rfloor$
        (lags up to 1~s).
  \item TRA was defined as:
        \begin{equation}
          \text{TRA}(x,y) = \sqrt{
            \frac{1}{\det(\boldsymbol{\Sigma})} \cdot
            \frac{1}{L} \sum_{\ell=1}^{L}
            \left( \frac{[y \star x]_\ell - [x \star y]_\ell}{N - \ell} \right)^2
          },
          \label{eq:tra}
        \end{equation}
        where $\boldsymbol{\Sigma}$ is the $2 \times 2$ covariance matrix of
        the normalised envelopes and $N$ is the signal length. Pairs for which
        $\det(\boldsymbol{\Sigma}) \leq 10^{-14}$ were excluded.
\end{enumerate}

Electrode-level TRA was defined as the mean TRA across all valid pairs
involving that electrode.

\subsection{Statistical analysis}

\subsubsection{Spearman rank correlations}

Associations between electrode-level waveform features and connectivity
measures were quantified using Spearman's rank correlation coefficient
$\rho_s$~\citep{spearman1904proof}:
%
\begin{equation}
  \rho_s = 1 - \frac{6 \sum d_i^2}{n(n^2-1)},
\end{equation}
%
where $d_i$ is the difference between the ranks of the $i$-th observation
on the two variables and $n$ is the number of observations. Spearman's $\rho$
was chosen because neither waveform features nor connectivity measures are
normally distributed, and the relationship is not assumed linear. Correlations
were computed (i)~pooled across all $n = 1{,}009$ electrodes (all subjects
combined), and (ii)~within each participant separately (using that
participant's electrodes only).

\subsubsection{Group-level inference: Wilcoxon signed-rank test}

The eight within-subject $\rho_s$ values for each feature--connectivity pair
were tested against zero using the Wilcoxon signed-rank test
\citep{wilcoxon1945individual}, a non-parametric test that evaluates whether
the median of a set of signed differences is consistently positive or negative.
This approach treats each participant as an independent observation and is
robust to non-normality of the $\rho$ distribution. With $n = 8$, the minimum
achievable two-tailed $p$-value is approximately 0.008.

\subsubsection{False discovery rate correction}

Multiple comparisons across nine hypothesis tests (3~waveform
features~$\times$~3~connectivity measures) were controlled using the
Benjamini--Hochberg procedure~\citep{benjamini1995controlling} at a nominal
\fdr level of 0.05. FDR correction was applied jointly across all nine tests
within a given analysis (pooled or \lme); uncorrected $\rho$ and $p$-values
are shown in the figures for completeness.

\subsubsection{Linear mixed-effects models}

To isolate the contribution of waveform shape features to connectivity
independent of anatomical and experimental confounds, we fitted three
linear mixed-effects models (one per connectivity outcome) using
\texttt{statsmodels.MixedLM}:
%
\begin{equation}
  y_z^{(e)} = \beta_0 + \beta_1 \cdot \text{IF}_z^{(e)}
             + \beta_2 \cdot \text{asc2desc}_z^{(e)}
             + \beta_3 \cdot \text{peak2trough}_z^{(e)}
             + \gamma_{\text{region}(e)}
             + \gamma_{\text{freq}(e)}
             + \gamma_{\text{stim}(e)}
             + u_{\text{subj}(e)}
             + \varepsilon^{(e)},
  \label{eq:lme}
\end{equation}
%
where $y_z$ is the z-scored connectivity outcome, subscript~$z$ denotes
z-scoring, $e$ indexes electrodes, $\gamma_{\cdot}$ are fixed-effect
contrasts for brain region (seven-level categorical; reference level:
frontal; amygdala merged with ``other'' due to $n = 2$), stimulation
frequency (1~Hz vs 50~Hz), and stimulated-channel status (binary).
$u_{\text{subj}(e)} \sim \mathcal{N}(0, \tau^2)$ is a random intercept
for the participant to which electrode $e$ belongs, accommodating
between-participant baseline differences in connectivity. $\varepsilon^{(e)}
\sim \mathcal{N}(0, \sigma^2)$ is the residual error.

Models were estimated by restricted maximum likelihood (REML) using the
Powell optimiser; the lbfgs optimiser was found to produce boundary
random-effect variance estimates ($\hat{\tau}^2 = 0$) and was therefore
avoided. Where $\hat{\tau}^2 < 10^{-6}$ after Powell optimisation,
the model was re-estimated by maximum likelihood. Standardised
coefficients ($\hat{\beta}$) and their 95\% confidence intervals are
reported; \fdr correction (Benjamini--Hochberg) was applied across all
nine waveform-feature fixed effects jointly (Table~\ref{tab:lme}).

\subsection{Waveform normalisation for visualisation}

For the pre- vs post-stimulation waveform comparison (Figure~\ref{fig:fig2}),
each electrode's alpha waveform was represented in a phase-aligned frame.
For each accepted cycle, the instantaneous amplitude $A(t)$ was used to
normalise the IMF~4 signal by its minimum and maximum, yielding a waveform
in the range $[-1, 1]$. These were then interpolated to a common 49-point
phase axis ($0 \to 1$ corresponding to ascending zero-crossing to ascending
zero-crossing), and averaged across all cycles within an electrode. Grand
averages across electrodes were computed by taking the mean and standard
error across electrodes, giving each electrode equal weight regardless of its
cycle count.

\subsection{Software and reproducibility}

All analyses were implemented in Python~3. Decomposition used the
\texttt{emd} package~\citep{quinn2021within}; connectivity measures used
\texttt{mne-connectivity} (ImCoh) and \texttt{statsmodels} (VAR/gPDC and
\lme); signal processing used \texttt{scipy}; statistical tests used
\texttt{scipy.stats} and \texttt{statsmodels.stats}; figures used
\texttt{matplotlib} with Arial fonts at 300~dpi (PDF fonttype 42).
All analysis code, configuration files, and figure scripts are available
at [repository URL]. Random seeds were fixed (\texttt{seed = 20260612}).

% ─────────────────────────────────────────────────────────────────────────────
\section*{Data availability}

The anonymised electrode-level summary data supporting the statistical
analyses in this paper are available from the corresponding author on
reasonable request. Raw \seeg recordings cannot be shared owing to patient
confidentiality constraints.

\section*{Code availability}

All analysis scripts are available at [GitHub/OSF repository URL].

\section*{Author contributions}

[To be completed.]

\section*{Competing interests}

The authors declare no competing interests.

\section*{Acknowledgements}

[To be completed.]

% ─────────────────────────────────────────────────────────────────────────────
\bibliographystyle{naturemag}
\bibliography{references}

% ─────────────────────────────────────────────────────────────────────────────
%% TABLES
\clearpage
% Table 1 — Participant and recording demographics
\begin{table}[h!]
\caption{\textbf{Participant and recording characteristics.}
sEEG = stereoelectroencephalography; CAR = common-average reference;
n\textsubscript{elec} = number of electrodes passing quality control in the
post-stimulation epoch (CAR montage); IF = instantaneous frequency of IMF\,4.
}
\label{tab:demographics}
\centering
\small
\begin{tabular}{lcccccc}
\toprule
\textbf{Participant} & \textbf{Stim freq (Hz)} & \textbf{n\textsubscript{elec}} &
\textbf{Mean IF (Hz)} & \textbf{Mean asc2desc} & \textbf{Mean peak2trough} &
\textbf{Region (majority)} \\
\midrule
Subject01         & 1, 50 &  65 & 9.74 $\pm$ 0.61 & 0.504 $\pm$ 0.005 & 0.490 $\pm$ 0.003 & Frontal \\
Subject04         & 1, 50 &  57 & 9.91 $\pm$ 0.82 & 0.506 $\pm$ 0.006 & 0.490 $\pm$ 0.003 & Frontal \\
Subject07         & 1, 50 &  16 & 9.98 $\pm$ 0.73 & 0.508 $\pm$ 0.005 & 0.489 $\pm$ 0.002 & Frontal \\
HEJ\_Subject12   & 1, 50 & 181 & 9.87 $\pm$ 0.76 & 0.506 $\pm$ 0.006 & 0.490 $\pm$ 0.003 & Temporal \\
HEJ\_Subject13   & 1, 50 & 201 & 9.96 $\pm$ 0.74 & 0.506 $\pm$ 0.005 & 0.490 $\pm$ 0.003 & Frontal \\
HEJ\_Subject14   & 1, 50 & 158 & 9.85 $\pm$ 0.78 & 0.506 $\pm$ 0.006 & 0.491 $\pm$ 0.003 & Temporal \\
HEJ\_Subject15   & 1, 50 & 258 & 9.83 $\pm$ 0.75 & 0.505 $\pm$ 0.005 & 0.490 $\pm$ 0.003 & Frontal \\
HEJ\_Subject16   & 1    &  73 & 9.90 $\pm$ 0.77 & 0.506 $\pm$ 0.006 & 0.490 $\pm$ 0.003 & Temporal \\
\midrule
\textbf{Total}   & --   & \textbf{1009} & \textbf{9.88 $\pm$ 0.76} & \textbf{0.506 $\pm$ 0.006} & \textbf{0.490 $\pm$ 0.003} & Mixed \\
\bottomrule
\end{tabular}
\end{table}

\clearpage
% Table 2 — LME fixed-effect results
\begin{table}[h!]
\caption{\textbf{Linear mixed-effects model results for waveform shape predicting
connectivity.}
Each row corresponds to one waveform-feature fixed effect in a model predicting the
z-scored connectivity measure (one model per outcome). All continuous predictors and
outcomes were z-scored prior to fitting. Fixed-effect confounds included brain region
(seven-level categorical; reference = frontal), stimulation frequency (1~Hz vs 50~Hz),
and stimulated-channel status. A random intercept per participant was included.
$\hat{\beta}$ = standardised coefficient; 95\%~CI = 95\% confidence interval;
$p$ = likelihood-ratio $p$-value; $q$ = Benjamini--Hochberg FDR-corrected $p$-value
across all nine waveform-feature tests (3~features~$\times$~3~outcomes).
Bold rows: $q < 0.05$.
}
\label{tab:lme}
\centering
\small
\begin{tabular}{llccccc}
\toprule
\textbf{Outcome} & \textbf{Predictor} & $\hat{\beta}$ & \textbf{95\% CI} & $p$ & $q$ & \\
\midrule
\multirow{3}{*}{$|$ImCoh$|$}
 & \textbf{Instantaneous frequency} & \textbf{0.117} & \textbf{[0.046, 0.187]} & \textbf{0.0012} & \textbf{0.0026} & \textbf{**} \\
 & \textbf{Ascent/descent ratio}    & \textbf{0.129} & \textbf{[0.073, 0.185]} & \textbf{$<$0.001} & \textbf{$<$0.001} & \textbf{***} \\
 & Peak/trough ratio                & $-$0.039      & [$-$0.095, 0.018]      & 0.184  & 0.275 & ns \\
\midrule
\multirow{3}{*}{PDC asymmetry}
 & Instantaneous frequency          & 0.012         & [$-$0.063, 0.087]      & 0.752  & 0.752 & ns \\
 & \textbf{Ascent/descent ratio}    & \textbf{$-$0.097} & \textbf{[$-$0.158, $-$0.035]} & \textbf{0.0021} & \textbf{0.0037} & \textbf{**} \\
 & \textbf{Peak/trough ratio}       & \textbf{0.168} & \textbf{[0.113, 0.224]} & \textbf{$<$0.001} & \textbf{$<$0.001} & \textbf{***} \\
\midrule
\multirow{3}{*}{TRA}
 & Instantaneous frequency          & $-$0.043      & [$-$0.112, 0.027]      & 0.228  & 0.293 & ns \\
 & \textbf{Ascent/descent ratio}    & \textbf{0.094} & \textbf{[0.039, 0.149]} & \textbf{$<$0.001} & \textbf{0.0026} & \textbf{***} \\
 & Peak/trough ratio                & $-$0.010      & [$-$0.066, 0.046]      & 0.734  & 0.752 & ns \\
\bottomrule
\end{tabular}
\end{table}


% ─────────────────────────────────────────────────────────────────────────────
%% FIGURES
\clearpage

\begin{figure}[p]
  \centering
  \includegraphics[width=\textwidth]{figures/fig1_emd_pipeline.png}
  \caption{\textbf{EMD waveform analysis pipeline applied to a representative
  sEEG electrode.}
  \textbf{A,} Six seconds of raw sEEG signal (top, dark) decomposed into six
  intrinsic mode functions (IMF~1--6; grey). IMF~4, capturing the alpha
  oscillation (9.9~Hz), is highlighted in blue with background shading.
  Traces are independently normalised and stacked with uniform spacing; scale
  bar = 1~s.
  \textbf{B,} A single representative alpha cycle from IMF~4 (blue), with key
  landmarks annotated: peak (red circle), trough (gold circle), and descending
  zero-crossing (green circle). Background grey traces show adjacent cycles.
  The ascent/descent ratio (\texttt{asc2desc}; Eq.~\ref{eq:asc2desc}) and
  peak/trough ratio (\texttt{peak2trough}; Eq.~\ref{eq:peak2trough}) are
  indicated by horizontal arrows.
  \textbf{C,} Phase-aligned instantaneous-frequency (PA-IF) heatmap for all
  accepted post-stimulation cycles (colour scale: Hz). Each row is one cycle;
  the $x$-axis spans 0--360$^\circ$. The white dashed line shows the mean
  IF profile across cycles.
  \textbf{D,} Population distributions of the three waveform features across
  all $n = 1{,}009$ post-stimulation electrodes (CAR montage). Vertical dashed
  lines mark medians.
  }
  \label{fig:fig1}
\end{figure}

\begin{figure}[p]
  \centering
  \includegraphics[width=\textwidth]{figures/fig2_waveform_change.png}
  \caption{\textbf{Electrical stimulation selectively alters ascending-phase
  duration at stimulated electrodes.}
  \textbf{A,} Grand mean ($\pm$ SEM) normalised alpha waveform across all
  $n = 42$ stimulated electrodes (pooled across 8 participants), pre-stimulation
  (blue) and post-stimulation (red). Each electrode contributes its mean cycle
  waveform, interpolated to a 49-point phase axis. Post-stimulation waveforms
  show a clear prolongation of the ascending phase.
  \textbf{B,} As in \textbf{A} but for all $n = 1{,}248$ non-stimulated
  electrodes. Pre- and post-stimulation waveforms are virtually identical,
  confirming that observed changes at stimulated sites are not due to global
  temporal drift.
  \textbf{C,} Violin plots comparing $\Delta$\texttt{asc2desc} (post\,$-$\,pre)
  between non-stimulated and stimulated electrodes. The stimulated group shows
  a significant positive shift (Mann--Whitney $p = 5.1 \times 10^{-8}$,
  $^{***}$); the non-stimulated group is centred at zero.
  \textbf{D,} As in \textbf{C} for $\Delta$\texttt{peak2trough}, showing no
  significant group difference ($p = 0.17$, ns).
  }
  \label{fig:fig2}
\end{figure}

\begin{figure}[p]
  \centering
  \includegraphics[width=\textwidth]{figures/fig3_main_scatter.png}
  \caption{\textbf{Waveform shape predicts network connectivity across
  post-stimulation electrodes.}
  Each panel shows a scatter plot of one waveform feature ($x$-axis) against
  one connectivity measure ($y$-axis) for all $n = 1{,}009$ post-stimulation
  electrodes (CAR montage; one point per electrode). The red line is an
  OLS regression fit included as a visual aid; inference is based on Spearman
  $\rho$ (annotated in each panel). Rows: connectivity measure (ImCoh,
  PDC~asymmetry, TRA). Columns: waveform feature (instantaneous frequency,
  ascent/descent ratio, peak/trough ratio). Panels with $p < 0.05$ (uncorrected)
  are indicated by asterisks (* $p < 0.05$, ** $p < 0.01$, *** $p < 0.001$).
  }
  \label{fig:fig3}
\end{figure}

\begin{figure}[p]
  \centering
  \includegraphics[width=\textwidth]{figures/fig4_lme_coefficients.png}
  \caption{\textbf{Linear mixed-effects model confirms independent waveform
  shape effects after controlling for confounds.}
  Standardised fixed-effect coefficients ($\hat{\beta}$; Eq.~\ref{eq:lme})
  from separate \lme models predicting each connectivity outcome
  (ImCoh in blue, PDC~asymmetry in orange, TRA in green). Points show
  $\hat{\beta}$; horizontal bars show 95\% confidence intervals. Filled
  circles indicate tests surviving Benjamini--Hochberg \fdr correction
  ($q < 0.05$); open circles are non-significant. All continuous predictors
  and outcomes were z-scored, making coefficients directly comparable.
  See Table~\ref{tab:lme} for full numerical results.
  }
  \label{fig:fig4}
\end{figure}

\begin{figure}[p]
  \centering
  \includegraphics[width=\textwidth]{figures/fig5_lme_partial_regression.png}
  \caption{\textbf{Partial regression plots showing confound-adjusted
  waveform--connectivity relationships.}
  Each panel shows the partial residual scatter for one waveform feature
  (rows) against one connectivity measure (columns), after removing
  variance attributable to brain region, stimulation frequency, and
  stimulated-channel status. The $x$-axis is the z-scored waveform feature
  residualised from the \lme confounds; the $y$-axis is the \lme residual
  plus the predictor contribution ($y_{\text{resid}} + \hat{\beta} x$).
  Red line: OLS fit (slope equals $\hat{\beta}$ exactly in partial residual
  space). Annotated $\hat{\beta}$, $p$, and FDR-corrected $q$ values
  correspond to those in Table~\ref{tab:lme}.
  }
  \label{fig:fig5}
\end{figure}

\begin{figure}[p]
  \centering
  \includegraphics[width=0.95\textwidth]{figures/fig6_stim_vs_nonstim.png}
  \vspace{0.5cm}
  \includegraphics[width=0.95\textwidth]{figures/fig6b_1hz_vs_50hz.png}
  \caption{\textbf{Waveform--connectivity associations are consistent across
  stimulation status and protocol.}
  \textbf{Top row,} Instantaneous frequency vs each connectivity measure,
  shown separately for stimulated (red; $n = 28$) and non-stimulated
  (blue; $n = 981$) electrodes. Spearman $\rho$ and $p$-values for each
  subgroup are annotated. Non-stimulated electrodes show robust significant
  associations across all three connectivity measures.
  \textbf{Bottom row,} As above, stratified by stimulation frequency:
  1~Hz (blue; $n = 479$) vs 50~Hz (red; $n = 530$). Both protocols
  yield significant PDC and TRA associations, indicating that findings
  are not specific to one stimulation paradigm.
  }
  \label{fig:fig6}
\end{figure}

\begin{figure}[p]
  \centering
  \includegraphics[width=\textwidth]{figures/fig7_region.png}
  \caption{\textbf{Regional specificity of waveform--connectivity associations.}
  Spearman $\rho$ between instantaneous frequency and each connectivity
  measure (panels A--C), computed separately within each anatomical region
  (AAL atlas). Bar length encodes $\rho$; electrode counts ($n$) and
  FDR-corrected significance markers (* $q < 0.05$) are annotated.
  Frontal and insula regions show the most consistent associations.
  }
  \label{fig:fig7}
\end{figure}

\end{document}
